% !TEX root = ../glossario.tex
% ---------------------------------------------------------------------
%  Lettera A del Glossario. Una voce per comando, in ordine alfabetico:
%
%     \termine{Voce}{Definizione: una o due frasi chiare.}
%
%  Nella definizione si possono usare \emph{corsivo}, \textbf{grassetto},
%  \url{https://...} e \glo{altro termine} (se anche quello è nel
%  Glossario). Niente tabelle, immagini o elenchi.
%  Per una lettera nuova: copiate questo file con il nome della lettera
%  (es. capitoli/b.tex) e cambiate la riga \letteraglossario.
% ---------------------------------------------------------------------
\phantomsection
\letteraglossario{A}

\termine{Angular}{Framework open source basato su TypeScript, mantenuto da Google, impiegato per lo sviluppo di applicazioni web lato client (front-end).}

\termine{API}{\emph{Application Programming Interface}: insieme di regole e di punti
  di accesso con cui un componente software mette dati e funzionalità
  a disposizione di altri programmi, che possono usarli senza
  conoscerne l'implementazione.}

\termine{Artefatto}{Oggetto tangibile o documento prodotto durante le attività di sviluppo del software (es. modelli, codice sorgente, manuali, file di configurazione, specifiche).}

\termine{AWS}{\emph{Amazon Web Services}: piattaforma di servizi di cloud computing offerta da Amazon che mette a disposizione potenza di calcolo, archiviazione e altre funzionalità on-demand.}
